% $Id:context-scala.tex  $
% !TEX root = main.tex

\section{\ac{COP} in Scala using Predicated Generic Functions}
\label{sec:context-scala}


%%
\subsection{Predicated Functions Based on the State Monad}
\label{sec:predicaote-monad}

In this work we implemented a light-weight \ac{COP} extension for Scala based on Predicated Generic Functions. Nonetheless, our implementation has all components to enable dynamic software adaptations, following the scoping and activation rules of \ac{COP} languages.
Recall from \fref{sec:preliminaries.cop}, that layers encapsulate \emph{partial} behaviors which are in scope only when they are determined to be by the language's activation mechanism.
In our implementation, we define layers using the \lstinline|PredicatedFunction| abstraction. A \lstinline|PredicatedFunction| is the combination of an activation condition and a \lstinline|Behavior|.
The activation condition is a predicate that depends on the state of the current execution context, while the \lstinline|Behavior| is a function that may manipulate said state and return some result.
A \lstinline|PredicatedFunction| is then a partial function that depends on a given state value (\ie context), whose \lstinline|Behavior| can only be applied if the activation predicate evaluates to \lstinline|true| for said state value.

A \lstinline|Behavior| is represented through the \lstinline|State| Monad: a function that takes some state and returns a product containing the new state and a result, as shown in \fref{lst:state-type}

\begin{scala}[label=lst:state-type, caption=A simplified definition of the State monad]
type State[S, A] = S => (S, A)
\end{scala}

The \lstinline|State| type has a \lstinline|Functor| instance that provides a function to map over the result type, as defined in \fref{lst:state-functor}.

\begin{scala}[label=lst:state-functor, caption=Functor instance for the State type]
given [S]: Functor[[X] =>> State[S, X]] with

  def map[A, B](fa: State[S, A], f: A => B): State[S, B] =
    state0 =>
      val (state1, a) = fa(state0)
      (state1, f(a))

\end{scala}

Additionally, the \lstinline|State| type has the \lstinline|Monad| instance shown in \fref{lst:state-monad}, specifying how to chain two state functions sequentially,
by threading the initial state through both state functions

\begin{scala}[label=lst:state-monad, caption=Monad instance for the State type]
given [S]: Monad[[X] =>> State[S, X]] with

  def flatMap[A, B](fa: State[S, A], f: A => State[S, B]): State[S, B] =
    state0 =>
      val (state1, a) = fa(state0)
      f(a)(state1)

\end{scala}

We saw in \fref{sec:preliminaries.cop} that multiple layers can be composed together to form new behavioral variations,
and that each layer within the composition can access the behavior from a previously active layer.
To achieve the same functionality with \lstinline|PredicatedFunctions|, we define a \lstinline|Behavior| to be a \lstinline|State| function that depends on a previously active computation,
as illustrated in \fref{lst:behavior}

\begin{scala}[label=lst:behavior, caption=The definition of a Behavior]
type Behavior[S, A] = A => State[S, A]
\end{scala}

Such representation enables a composable definition of \ac{COP}. \text{Layer} composition is translated to function composition,
which is easier to implement. Rather than having a special \lstinline|proceed| construct that would require meta-programming to implement,
we use Higher-Order Functions to express state manipulations that may depend on previous computations.

Finally, with these definitions in place, A \lstinline|PredicatedFunction| is defined in \fref{lst:predicated-function} as a partial \lstinline|Behavior|

\begin{scala}[label=lst:predicated-function, caption=The definition of a Predicated Function]
trait PredicatedFunction[S, A]:

  def isActive(state: S): Boolean

  def run: Behavior[S, A]

\end{scala}

We use the \lstinline|Monad| instance of the \lstinline|State| type to compose various \lstinline|Behaviors|.
There are two strategies to combine two \lstinline|PredicatedFunctions|.
The first strategy consist on simply running them one after the other and feeding the result of the
first into the second one. If only one of the two Behaviors is active for the current state, their respective result is returned, otherwise, if both are active
at the same time, both Behaviors are chained using \lstinline|flatMap|. If neither is active, then none of the Behaviors are executed, and the result of the
most recently active Behavior is returned

\begin{scala}[label=lst:predicated-function-andThen, caption=Sequential Composition of Predicated Functions]
infix def andThen[S, A](
    f1: PredicatedFunction[S, A],
    f2: PredicatedFunction[S, A]
) = new PredicatedFunction[S, A]:

  override def isActive(state: S): Boolean =
    f1.isActive(state) || f2.isActive(state)

  override val run: Behavior[S, A] = previous =>
    State.get.flatMap { state =>
      (f1.isActive(state), f2.isActive(state)) match
        case (true, false) => f1.run(previous)
        case (false, true) => f2.run(previous)
        case (true, true)  => f1.run(previous).flatMap(f2.run)
        case _             => passthrough(previous)
    }
\end{scala}

The following auxiliary functions (\fref{lst:predicated-function-andThen-aux}) were used to define the composition function in \fref{lst:predicated-function-andThen}:

\begin{scala}[label=lst:predicated-function-andThen-aux, caption=Auxiliary functions to implement Predicated Functions]

/**
* A state transition that does not modify the state, and returns the result of the previous computation
*/
def passthrough[S, A]: Behavior[S, A] =
  previous => state => (state, previous)

object State:
  // Get the current state
  def get[S]: State[S] = state => (state, state)
\end{scala}

The second strategy to compose two Predicated Functions is shown in \fref{lst:predicated-function-orElse}, and it consists on creating another function which
chooses the first active behavior. Instead of chaining both functions, only one will be chosen
and executed.

\begin{scala}[label=lst:predicated-function-orElse, caption=Alternating between Predicated Functions]
infix def orElse[S, A](f1: PredicatedFunction[S, A], f2: PredicatedFunction[S, A]) = new PredicatedFunction[S, A]:

  override def isActive(state: S): Boolean =
  f1.isActive(state) || f2.isActive(state)

  override val run: Behavior[S, A] = previous =>
  State.get.flatMap { state =>
    (f1.isActive(state), f2.isActive(state)) match
      case (true, _)     => f1.run(previous)
      case (false, true) => f2.run(previous)
      case _             => passthrough(previous)
  }
\end{scala}

Using these binary operations, we can easily create functions to compose a sequence of Predicated Functions:

\begin{scala}[label=lst:predicated-function-compose-multiple, caption=Composing multiple Predicated Functions]
object PredicatedFunction:

  def sequence[S, A](
      functions: NonEmptyList[PredicatedFunction[S, A]]
  ): PredicatedFunction[S, A] =
    functions.reduceLeft(_ andThen _)

  def oneOf[S, A](
      functions: NonEmptyList[PredicatedFunction[S, A]]
  ): PredicatedFunction[S, A] =
    functions.reduceLeft(_ orElse _)

\end{scala}

Finally, Predicated Functions can be created by extending the \lstinline|PredicatedFunction| trait as
shown in \fref{lst:predicated-function-creation-example}. This example shows three \lstinline|PredicatedFunction| that
depend on a context state of type \lstinline|Int|, and return a result of type \lstinline|String|.
Each function defines a \lstinline|Behavior| that may manipulate the state in some way and return a result.
For instance, \fref{ln:predicated-function-creation-example.fun2} shows that, when \lstinline|Fun2| is active, it adds 10 to the \lstinline|state| value, and it appends the \lstinline|"Hello from Fun2. "| string to
the result of the most recently active \lstinline|PredicatedFunction|.

\begin{scala}[label=lst:predicated-function-creation-example, caption=Predicated Function examples, escapechar=~]
object Fun1 extends PredicatedFunction[Int, String]:
  override def isActive(state: Int): Boolean =
    state < 0

  override def run: Behavior[Int, String] = previous => state =>
    (state, "Hello from Fun1. ")

object Fun2 extends PredicatedFunction[Int, String]:
  override def isActive(state: Int): Boolean =
    state > 0

  override def run: Behavior[Int, String] = previous => state =>
    (state + 10, previous + "Hello from Fun2. ") ~\label{ln:predicated-function-creation-example.fun2}~

object Fun3 extends PredicatedFunction[Int, String]:
  override def isActive(state: Int): Boolean =
    state > 10

  override def run: Behavior[Int, String] = previous => state =>
    (state - 2, previous + "Hello from Fun3. ")
\end{scala}

\fref{lst:predicated-function-composition-example} demonstrates how to combine the Behaviors defined in \fref{lst:predicated-function-creation-example}
using the two composition strategies: \lstinline|sequence| (\fref{ln:predicated-function-composition-example.fun4}) and \lstinline|oneOf| (\fref{ln:predicated-function-composition-example.fun5})

\begin{scala}[label=lst:predicated-function-composition-example, caption=Predicated Function composition examples, escapechar=~]
val initialState = 1
val initialResult = ""

val Fun4 = PredicatedFunction.sequence(Fun1, Fun2, Fun3) ~\label{ln:predicated-function-composition-example.fun4}~
val Fun5 = PredicatedFunction.oneOf(Fun1, Fun2, Fun3) ~\label{ln:predicated-function-composition-example.fun5}~

val run1 = Fun4.run(initialState)(initialResult) ~\label{ln:predicated-function-composition-example.run1}~
// run1 = (9, "Hello from Fun2. Hello from Fun3. ")
val run2 = Fun5.run(initialState)(initialResult) ~\label{ln:predicated-function-composition-example.run2}~
// run2 = (11, "Hello from Fun2. ")
\end{scala}

Given an initial state of 1, the result of the \lstinline|sequence| composition in \fref{ln:predicated-function-composition-example.run1} is \lstinline|(9, "Hello from Fun2. Hello from Fun3. ")|
because only \lstinline|Fun2| and \lstinline|Fun3| are active for that particular value. \lstinline|Fun2| first adds 10 to the initial state, which then activates the Behavior of \lstinline|Fun3| because its predicate requires that \lstinline|state > 10|.
The result of the \lstinline|oneOf| composition (\fref{ln:predicated-function-composition-example.run2}), on the other hand, is \lstinline|(11, "Hello from Fun2. ")| because only the first active function in the composition is executed.

Note that it is not possible to get into a situation where a Predicated Function does not have access to
a \lstinline|previous| computation. This is because their definition requires the programmer to always
give an initial state and an initial value of the result. In the case where no Functions are active,
the result will be equal to those initial values. As it will be shown in the next chapter, this simple way of
structuring programs matches well with the chosen architecture for adaptive systems.


%%
\subsection{Refined Predicated Functions}
\label{sec:refined-predicated-functions}

The \lstinline|PredicatedFunction| abstraction is enough to construct \ac{COP} programs, however it still has many of the
disadvantages explored in \fref{sec:problem-of-verifying-cop} when reasoning about the result of a functional composition. Take for instance the code in
\fref{lst:predicated-function-composition-example}. The value of \lstinline|run1| would be different if the programmer composed the functions in a different order;
for example, the composition \lstinline|sequence(Fun1, Fun3, Fun2)| yields \lstinline|(11, "Hello from Fun2. ")|, when applied to the pair of initial state and result \lstinline|(1, "")|.
Furthermore, the state manipulation of a function may contradict the activation predicate of the next function in the composition chain.
\fref{lst:predicated-function-contradictory-composition} shows an example where the first function manipulates the state in such a way that the next function will never be activated:


\begin{scala}[label=lst:predicated-function-contradictory-composition, caption=Contradictory composition example]

object Fun2 extends PredicatedFunction[Int, String]:
  override def isActive(state: Int): Boolean =
    state > 0

  override def run: Behavior[Int, String] = previous => state =>
    (state - 2 * state, previous + "Hello from Fun2. ")

object Fun3 extends PredicatedFunction[Int, String]:
  override def isActive(state: Int): Boolean =
    state > 10

  override def run: Behavior[Int, String] = previous => state =>
    (state - 2, previous + "Hello from Fun3. ")

// Fun3 is never active whenever Fun2 is, for all possible values of the input state
val Fun6 = PredicatedFunction.sequence(Fun2, Fun3)
\end{scala}

Whether the state manipulations and composition order are coherent or not, depends on the requirements of the program.
It is the responsibility of the programmer to enforce and validate that these requirements are fulfilled. Though this an ubiquitous problem for any
program (irrespective of its programming paradigm), \ac{COP} programs pose additional challenges when verifying their behavior due to the notion of activation.
The programmer can (and should) define tests to validate that the results for these compositions make sense from a requirements standpoint; however, as stated
in \fref{sec:problem-of-verifying-cop}, a comprehensive validation of a \ac{COP} program can be complex and expensive due to the combinatorial explosion of possible states and layer interactions.

Our proposal takes on a different approach to verify such programs; one that enables the programmer to express properties about activations and interactions, without actually having to incur in the cost of exploring the vast space of possible states.
Using Refinement Types, one can embed Classical Logic statements into a type and automatically verify their validity.
The main modification that is done to the \lstinline|PredicatedFunction|, is to simply add the activation predicate as a precondition to the \lstinline|run| function, as show in \fref{lst:predicated-function-run-precondition}.

\begin{scala}[label=lst:predicated-function-run-precondition, caption=Refinement of the input state type]
abstract class PredicatedFunction[S, A]:

  def isActive(state: S): Boolean

  def run(previous: A, state: S): (S, A) =
    require(this.isActive(state))
    (??? : (S, A))
\end{scala}

Such change refines the type of the \lstinline|run| function to:

$$
run : A \rightarrow S\{state: S  \: | \: isActive(state) \rightarrow (S, A)\}
$$

Adding the pre-condition to the \lstinline|run| function's \lstinline|state| argument means that it can be only applied if the caller provides an \emph{evidence} that the \lstinline|isActive| predicate holds for the given \lstinline|state| input.
Such specification also allows the programmer to bring the assumptions from the \lstinline|isActive| predicate into scope, so that they can be used when defining the program's logic. \fref{lst:predicated-function-run-precondition-example} shows
a simple example of how an optional value can be safely accessed from within the \lstinline|run| function's scope, due to the \lstinline|isActive| precondition.
The \lstinline|state| function parameter is refined to \lstinline|Option[BigInt]{v : v.isInstanceOf[Some] && v.get > 10}|, guaranteeing that \lstinline|state.get| will never throw an exception as it should always be an instance of the \lstinline|Some| type.


\begin{scala}[label=lst:predicated-function-run-precondition-example, caption=Refinement Typed-Checked example of the run function]
object RefinedFun1 extends PredicatedFunction[Option[BigInt], String]:

  override def isActive(state: Option[BigInt]): Boolean =
    state match
      case Some(x) => x > 10
      case None    => false

  override def run(
      previous: String,
      state: Option[BigInt]
  ): (Option[BigInt], String) =
    // The .get call is safe and verified at compile-time.
    // The isActive precondition states that this function
    // can only be called when the optional value is defined
    (state.get - 9, "Hello from Fun1. ")
\end{scala}

From this simple restriction, more complex properties can be defined between two or more functions that interact with each other.
It is up to the programmer to define these properties per the program's requirements. For example, it is particularly useful
to ensure that the activation and execution of a predicated function does not deactivate the next function in the composition chain.
This can be formalized in Stainless as a predicate stating that the activation of the first function implies the activation of the second one.
It is not enough, however, to establish a relation between the activation predicates without taking into account the actual logic of the functions.
Recall from previous examples that each function can manipulate the state, and it may do so in such a way that the next function becomes deactivated for certain state values.
A proper formulation should then assert that the activation and state manipulation of the first function, implies that next one will always be active.
The next abstraction will aid in formalizing such property:

\begin{scala}[label=lst:predicated-function-chainable, caption=Chainable type for expressing and verifying dependencies between Predicated Functions, escapechar=~]
abstract class Chainable[S, A]:

  def first: PredicatedFunction[S, A]
  def second: PredicatedFunction[S, A]

  @law
  def isChainable(previous: A, state: S): Boolean =
    first.isActive(state) ==> {
      val (newState, _) = first.run(previous, state) ~\label{ln:predicated-function-chainable.first.run}~
      second.isActive(newState)
    }
\end{scala}

An instance of the \lstinline|Chainable| type serves as an evidence that, for the given \lstinline|first| and \lstinline|second| predicated functions,
the \lstinline|isChainable| predicate holds for all possible values of \lstinline|previous: A| and \lstinline|state: S|. The property states that the
activation condition of the \lstinline|first| function implies that, when applying it to the previous result and state, the \lstinline|second| function will be active for the new returned state.
The type of the \lstinline|isChainable| property is shown in Equation \ref{eqn:isChainable}.

\begin{equation}
\label{eqn:isChainable}
\begin{split}
isChainable &: (previous: A) \\
            &\rightarrow (state: S) \\
            &\rightarrow unit\{f1.isActive(state) \\
            &\Rightarrow f2.isActive(f1.run(previous, state).first)\}
\end{split}
\end{equation}


Note that the \lstinline|first.run| call on \fref{ln:predicated-function-chainable.first.run} must necessarily happen inside the right hand side block of the implication (\lstinline|==>|) because the \lstinline|run| function's precondition states that the \lstinline|isActive| predicate
should hold for the given state input.

The next code snippet shows two functions that satisfy this implication for all possible state values.

\begin{scala}[label=lst:predicated-function-chainable-example, caption=Chainable type usage example]
final class RefinedFun1 extends PredicatedFunction[Option[BigInt], String]:

  override def isActive(state: Option[BigInt]): Boolean =
    state match
      case Some(x) => x > 10
      case _       => false

  override def run(
      previous: String,
      state: Option[BigInt]
  ): (Option[BigInt], String) =
    (state.get - 9, "Hello from Fun1. ")

final class RefinedFun2 extends PredicatedFunction[Option[BigInt], String]:
  override def isActive(state: Option[BigInt]): Boolean =
    state match
      case Some(x) => x > 1
      case _       => false

  override def run(
      previous: String,
      state: Option[BigInt]
  ): (Option[BigInt], String) =
    (state.get + 1, "Hello from Fun2. ")

// Stainless automatically verifies that the isChainable law holds for the given functions
val chain = new Chainable[Option[BigInt], String]:
  def first: PredicatedFunction[Option[BigInt], String] = new RefinedFun1
  def second: PredicatedFunction[Option[BigInt], String] = new RefinedFun2
\end{scala}

When instantiating a new \lstinline|Chainable| instance with two concrete Predicated Functions,
Stainless will automatically verify that the \lstinline|isChainable| law holds for those particular functions.
If, for instance, \lstinline|RefinedFun1| subtracted 10 instead of 9 to the input state (\lstinline|state.get - 10|), then
the programmer would get a compilation error because there exists a \lstinline|state| value for which the \lstinline|isChainable| property does not hold;
the SMT solver will also return an example of a value that makes the predicate false (in this case \lstinline|state =  11|)


\endinput
